\documentclass[10pt,a4paper]{amsart}
\usepackage{amsmath,amssymb,amsthm,graphicx}
\usepackage{verbatim}
\newenvironment{reference}{\comment}{\endcomment}
%\newenvironment{reference}{}{}
\newenvironment{slogan}{\comment}{\endcomment}
\newenvironment{history}{\comment}{\endcomment}
% For commutative diagrams we use Xy-pic
\usepackage[all]{xy}

% We use 2cell for 2-commutative diagrams.
\xyoption{2cell}
% Theorem environments.
%
\theoremstyle{plain}
\newtheorem{theorem}[subsection]{Theorem}
\newtheorem{proposition}[subsection]{Proposition}
\newtheorem{lemma}[subsection]{Lemma}

\theoremstyle{definition}
\newtheorem{definition}[subsection]{Definition}
\newtheorem{example}[subsection]{Example}
\newtheorem{exercise}[subsection]{Exercise}
\newtheorem{situation}[subsection]{Situation}

\theoremstyle{remark}
\newtheorem{remark}[subsection]{Remark}
\newtheorem{remarks}[subsection]{Remarks}

\numberwithin{equation}{subsection}

% Macros
%
\def\lim{\mathop{\mathrm{lim}}\nolimits}
\def\colim{\mathop{\mathrm{colim}}\nolimits}
\def\Spec{\mathop{\mathrm{Spec}}}
\def\Hom{\mathop{\mathrm{Hom}}\nolimits}
\def\Ext{\mathop{\mathrm{Ext}}\nolimits}
\def\SheafHom{\mathop{\mathcal{H}\!\mathit{om}}\nolimits}
\def\SheafExt{\mathop{\mathcal{E}\!\mathit{xt}}\nolimits}
\def\Sch{\mathit{Sch}}
\def\Mor{\mathop{\mathrm{Mor}}\nolimits}
\def\Ob{\mathop{\mathrm{Ob}}\nolimits}
\def\Sh{\mathop{\mathit{Sh}}\nolimits}
\def\NL{\mathop{N\!L}\nolimits}
\def\CH{\mathop{\mathrm{CH}}\nolimits}
\def\proetale{{pro\text{-}\acute{e}tale}}
\def\etale{{\acute{e}tale}}
\def\QCoh{\mathit{QCoh}}
\def\Ker{\mathop{\mathrm{Ker}}}
\def\Im{\mathop{\mathrm{Im}}}
\def\Coker{\mathop{\mathrm{Coker}}}
\def\Coim{\mathop{\mathrm{Coim}}}

% Boxtimes
%
\DeclareMathSymbol{\boxtimes}{\mathbin}{AMSa}{"02}

%
% Macros for moduli stacks/spaces
%
\def\QCohstack{\mathcal{QC}\!\mathit{oh}}
\def\Cohstack{\mathcal{C}\!\mathit{oh}}
\def\Spacesstack{\mathcal{S}\!\mathit{paces}}
\def\Quotfunctor{\mathrm{Quot}}
\def\Hilbfunctor{\mathrm{Hilb}}
\def\Curvesstack{\mathcal{C}\!\mathit{urves}}
\def\Polarizedstack{\mathcal{P}\!\mathit{olarized}}
\def\Complexesstack{\mathcal{C}\!\mathit{omplexes}}
% \Pic is the operator that assigns to X its picard group, usage \Pic(X)
% \Picardstack_{X/B} denotes the Picard stack of X over B
% \Picardfunctor_{X/B} denotes the Picard functor of X over B
\def\Pic{\mathop{\mathrm{Pic}}\nolimits}
\def\Picardstack{\mathcal{P}\!\mathit{ic}}
\def\Picardfunctor{\mathrm{Pic}}
\def\Deformationcategory{\mathcal{D}\!\mathit{ef}}
\usepackage{cleveref}
\crefname{originallemma}{original lemma}{original lemmas}
\Crefname{originallemma}{Original lemma}{Original lemmas}
\makeatletter
\newlabel{lemma-DM}{{lemma-DM}{source}{}{sourcepointer-1}{}}
\newlabel{lemma-DM@cref}{{[originallemma][0][]lemma-DM}{[1][0][]source}{}{}{}}
\makeatother
\crefname{originallemma}{original lemma}{original lemmas}
\Crefname{originallemma}{Original lemma}{Original lemmas}
\makeatletter
\newlabel{lemma-compose-after-separated}{{lemma-compose-after-separated}{source}{}{sourcepointer-2}{}}
\newlabel{lemma-compose-after-separated@cref}{{[originallemma][0][]lemma-compose-after-separated}{[1][0][]source}{}{}{}}
\makeatother
\crefname{originalreference}{original reference}{original references}
\Crefname{originalreference}{Original reference}{Original references}
\makeatletter
\newlabel{spaces-morphisms-lemma-permanence-unramified}{{spaces-morphisms-lemma-permanence-unramified}{source}{}{sourcepointer-3}{}}
\newlabel{spaces-morphisms-lemma-permanence-unramified@cref}{{[originalreference][0][]spaces-morphisms-lemma-permanence-unramified}{[1][0][]source}{}{}{}}
\makeatother
\title{Unramifiedness is preserved by base change and composition}
\author{The Stacks Project Authors}
\date{stacks-project:stacks-morphisms.tex:local-snapshot-20261012; GNU FDL 1.2 or later}
\begin{document}
\maketitle
\noindent\textit{Editorial note.} All delivered mathematical text is the author's own, copied from the pinned Stacks Project chapter source sources/stacks-algebraic-stacks-morphisms/pinned-original-source.tex; no statement or proof step was rewritten, reordered or removed.\par\medskip
\noindent\textit{Editorial note.} Preamble provenance: the Stacks Project keeps one single shared preamble for every chapter; this card takes it directly from the staged sources/stacks-algebraic-stacks-morphisms/preamble.tex as three verbatim spans (lines 8-12, 14-18 and 154-217), so no derived file was written and nothing was added or redefined. The wrapper class and the base packages are supplied by the assembler.\par\medskip
\noindent\textit{Editorial note.} Bibliography provenance: neither the counted statement nor its proof contains a citation (this lane rescans the whole span before writing the card), so the card carries no bibliography and no entry was reconstructed.\par\medskip
\noindent\textit{Editorial note.} Selection provenance: this card was opened under the file-based control contract (CONTROL.md), with the mechanical pre-card gate run on its statement and proof interval before the card was written.\par\medskip
\noindent\textit{Allowed standard prerequisites (editorial).}\begin{itemize}
\item The characterization of unramified morphisms through the diagonal being locally of finite type and quasi-separated, or through the corresponding inertia condition. Used inside the counted proof (source lines 8429-8433). Reduces unramifiedness to a property of the diagonal, which behaves well under base change and composition.
\item The corresponding permanence statement for the diagonal of a composition of morphisms of algebraic spaces. Invoked inside the counted proof (source lines 8436-8440) as a cross-chapter pointer. Supplies the composition step after base change to a presentation.
\end{itemize}\medskip
\ifdefined\hypertarget\hypertarget{sourcepointer-1}{}\fi
\noindent\textit{Source reference.} Original lemma lemma-DM: sources/stacks-algebraic-stacks-morphisms/pinned-original-source.tex, source lines 7800--7800. Navigation pointer only: the counted text refers to the Stacks item lemma-DM, which belongs to another chapter of the project and is not part of this card. The pointer carries the target's own label as its number because this chapter's staged snapshot does not contain the other chapter's tag table; its content is not reproduced here.\par\medskip
\ifdefined\hypertarget\hypertarget{sourcepointer-2}{}\fi
\noindent\textit{Source reference.} Original lemma lemma-compose-after-separated: sources/stacks-algebraic-stacks-morphisms/pinned-original-source.tex, source lines 1089--1089. Navigation pointer only: the counted text refers to the Stacks item lemma-compose-after-separated, which belongs to another chapter of the project and is not part of this card. The pointer carries the target's own label as its number because this chapter's staged snapshot does not contain the other chapter's tag table; its content is not reproduced here.\par\medskip
\ifdefined\hypertarget\hypertarget{sourcepointer-3}{}\fi
\noindent\textit{Source reference.} Original reference spaces-morphisms-lemma-permanence-unramified: sources/stacks-algebraic-stacks-morphisms/pinned-original-source.tex, source lines 968--968. Navigation pointer only: the counted text refers to the Stacks item spaces-morphisms-lemma-permanence-unramified, which belongs to another chapter of the project and is not part of this card. The pointer carries the target's own label as its number because this chapter's staged snapshot does not contain the other chapter's tag table; its content is not reproduced here.\par\medskip
\begin{lemma}
\label{lemma-permanence-unramified}
Let $\mathcal{X} \to \mathcal{Y} \to \mathcal{Z}$ be
morphisms of algebraic stacks.
If $\mathcal{X} \to \mathcal{Z}$ is unramified and
$\mathcal{Y} \to \mathcal{Z}$ is DM, then
$\mathcal{X} \to \mathcal{Y}$ is unramified.
\end{lemma}

\begin{proof}
Assume $\mathcal{X} \to \mathcal{Z}$ is unramified. By
Lemma \ref{lemma-compose-after-separated} the morphism
$\mathcal{X} \to \mathcal{Y}$ is DM. Choose a commutative diagram
$$
\xymatrix{
U \ar[d] \ar[r] & V \ar[d] \ar[r] & W \ar[d] \\
\mathcal{X} \ar[r] & \mathcal{Y} \ar[r] & \mathcal{Z}
}
$$
with $U, V, W$ algebraic spaces,
with $W \to \mathcal{Z}$ surjective smooth,
$V \to \mathcal{Y} \times_\mathcal{Z} W$ surjective \'etale, and
$U \to \mathcal{X} \times_\mathcal{Y} V$ surjective \'etale
(see Lemma \ref{lemma-DM}). Then also
$U \to \mathcal{X} \times_\mathcal{Z} W$ is surjective and \'etale.
Hence we know that $U \to W$ is unramified and we have to show that
$U \to V$ is unramified. This follows from
Morphisms of Spaces, Lemma \ref{spaces-morphisms-lemma-permanence-unramified}.
\end{proof}

\end{document}
